\usepackage{fontspec}
\usepackage{polyglossia}
\setdefaultlanguage[numerals=eastern]{arabic}
\setotherlanguage{english}
\newfontfamily\arabicfont[Script=Arabic]{Arial}
\setmainfont{Arial}

%==================================================
% دمج عدة صفوف في خلية واحدة داخل الجداول
%==================================================
\usepackage{multirow}
%==================================================
% إدراج الصور والتحكم في حجمها وتدويرها
%==================================================
\usepackage{graphicx}

%==================================================
% دعم الخيار [H] لتثبيت الجداول والأشكال في مكانها
%==================================================
\usepackage{float}

%==================================================
% ضبط هوامش الصفحة وأبعاد الورقة
%==================================================
\usepackage{geometry}

\geometry{
    a4paper,
    landscape,
    left=1cm,
    right=1cm,
    bottom=1cm,
    top=0.5cm,
    includehead,
    headheight=35pt,
    headsep=0pt,
    footskip=25pt
}

%==================================================
% التحكم في كابشن الجداول وزيادة الارتفاع عن الجدول
%==================================================
\usepackage{caption}
\captionsetup[table]{
    font={Large,bf},
    skip=5pt
}


%==================================================
% Header & Footer
%==================================================

\usepackage{fancyhdr}

\setlength{\headheight}{20pt}
\setlength{\headsep}{12pt}
\setlength{\footskip}{25pt}

%==================================================
% Section - Subsection بدون ترقيم
%==================================================

\makeatletter

\newcommand{\CurrentSection}{}
\newcommand{\CurrentSubsection}{}

\renewcommand{\sectionmark}[1]{%
    \gdef\CurrentSection{#1}%
    \gdef\CurrentSubsection{}%
    \markright{\CurrentSection}%
}

\renewcommand{\subsectionmark}[1]{%
    \gdef\CurrentSubsection{#1}%
    \markright{\CurrentSection\ - \CurrentSubsection}%
}

\makeatother

%==================================================
% الرأس والتذييل
%==================================================

\pagestyle{fancy}
\fancyhf{}

% الرأس - يسار
\fancyhead[L]{%
    \large\bfseries\nouppercase{\rightmark}
}

% الرأس - يمين
\fancyhead[R]{%
    \large\bfseries
    المديرية العامة للتربية في محافظة النجف الاشرف/ قسم التخطيط التربوي/ شعبة الاحصاء
}

% التذييل
\fancyfoot[L]{%
    \Large\bfseries\thepage
}

\fancyfoot[C]{}

\fancyfoot[R]{\large\bfseries الإحصاء السنوي للعام الدراسي ٢٠٢٧/٢٠٢٦}

% خطوط الرأس والتذييل
\renewcommand{\headrulewidth}{2pt}
\renewcommand{\footrulewidth}{2pt}

%==================================================
% الصفحة الأولى (plain)
%==================================================

\fancypagestyle{plain}{%
    \fancyhf{}

    % الرأس - يسار
    \fancyhead[L]{%
        \large\bfseries\nouppercase{\rightmark}
    }

    % الرأس - يمين
    \fancyhead[R]{%
        \large\bfseries
        المديرية العامة للتربية في محافظة النجف الاشرف/قسم التخطيط التربوي/شعبة الاحصاء
    }

    % التذييل
    \fancyfoot[L]{%
        \Large\bfseries\thepage
    }

    \fancyfoot[C]{}

    \fancyfoot[R]{\large\bfseries الإحصاء السنوي للعام الدراسي ٢٠٢٧/٢٠٢٦}

    \renewcommand{\headrulewidth}{2pt}
    \renewcommand{\footrulewidth}{2pt}
}

%==================================================
% تطبيق التنسيق على أول صفحة أيضاً
%==================================================

\AtBeginDocument{%
    \pagestyle{fancy}
    \thispagestyle{fancy}
}

%==================================================
% تغيير حجم السكشن
%==================================================
\usepackage{titlesec}
\titleformat*{\section}{\Large\bfseries\color{blue}}
\titleformat*{\subsection}{\Large\bfseries\color{blue}}
\titleformat*{\subsubsection}{\Large\bfseries\color{blue}}

%==================================================
% ادراج ملفات بي دي اف
%==================================================
\usepackage{pdfpages}

%==================================================
% التحكم في ارتفاع الجداول
%==================================================
\renewcommand{\arraystretch}{1}

%==================================================
% الحزم الخاصة بالجداول
%==================================================
\usepackage{color}
\usepackage{rotating}
\usepackage{tabularray}
\definecolor{Silver}{rgb}{0.752,0.752,0.752}